\section{ชุดข้อมูลที่มีอยู่และช่องว่าง}

งานวิจัยด้านการตรวจสอบสภาพถนนด้วยเซ็นเซอร์สมาร์ตโฟนส่วนใหญ่สร้างชุดข้อมูลเฉพาะสำหรับงานวิจัยของตนเอง และไม่ได้เผยแพร่ชุดข้อมูลดังกล่าวสู่สาธารณะอย่างต่อเนื่อง ตารางที่ 2.3 แสดงชุดข้อมูลและงานวิจัยที่เกี่ยวข้องกับการเก็บข้อมูลจากสมาร์ตโฟนเพื่อประเมินสภาพถนน ชุดข้อมูลเหล่านี้ส่วนใหญ่เก็บข้อมูลจากหน่วยวัดความเฉื่อยเป็นหลัก เช่น มาตรวัดความเร่งและมาตรวัดการหมุน ร่วมกับข้อมูลตำแหน่งหรือความเร็วในบางงาน อย่างไรก็ตาม ยังไม่พบชุดข้อมูลสาธารณะที่เก็บจากบริบทถนนเมืองไทยโดยตรง และยังพบข้อจำกัดด้านการเผยแพร่ข้อมูล กระบวนการกำหนดป้ายกำกับ และความสามารถในการทำซ้ำของกระบวนการทดลอง

\begingroup
\fontsize{11pt}{13pt}\selectfont
\setlength{\tabcolsep}{2pt}
\renewcommand{\arraystretch}{1.15}
\setlength{\LTleft}{0pt}
\setlength{\LTright}{0pt}
\begin{longtable}{>{\raggedright\arraybackslash}p{0.15\textwidth}>{\raggedright\arraybackslash}p{0.10\textwidth}>{\raggedright\arraybackslash}p{0.22\textwidth}>{\raggedright\arraybackslash}p{0.29\textwidth}>{\raggedright\arraybackslash}p{0.15\textwidth}}
\caption{สรุปงานวิจัยด้านการเก็บข้อมูล IMU เพื่อประเมินสภาพถนน}\label{tab:imu-road-datasets}\\
\toprule
\textbf{งานวิจัย / ชุดข้อมูล} & \textbf{ประเทศ} & \textbf{ประเภทของป้ายกำกับ} & \textbf{ลักษณะข้อมูล} & \textbf{สถานะการเผยแพร่} \\
\midrule
\endfirsthead
\multicolumn{5}{l}{\textit{ตารางที่ \thetable\ (ต่อ)}}\\
\toprule
\textbf{งานวิจัย / ชุดข้อมูล} & \textbf{ประเทศ} & \textbf{ประเภทของป้ายกำกับ} & \textbf{ลักษณะข้อมูล} & \textbf{สถานะการเผยแพร่} \\
\midrule
\endhead
\bottomrule
\endfoot
\textcite{gonzalez2017} & เม็กซิโก & 5 ประเภท ได้แก่ หลุม ลูกระนาดโลหะ ลูกระนาดยางมะตอย ถนนสึกหรอ และถนนปกติ & มาตรวัดความเร่งแกนตั้ง เก็บจากยานพาหนะ 12 คัน และตำแหน่งการวางสมาร์ตโฟนหลายรูปแบบ & บทความเดิมระบุว่าเผยแพร่สาธารณะ แต่รายงานล่าสุดระบุว่าไม่สามารถเข้าถึงได้แล้ว \\
\textcite{khandakar2025} & บังกลาเทศ & เหตุการณ์ถนนผิดปกติ เช่น ลูกระนาดและหลุม รวมถึงข้อมูลพฤติกรรมการขับขี่ & ข้อมูลหลายเซ็นเซอร์ ได้แก่ มาตรวัดความเร่ง มาตรวัดการหมุน สนามแม่เหล็ก ระบบระบุตำแหน่งบนพื้นโลก แรงโน้มถ่วง ทิศทาง ความเร่งรวม และข้อมูลเซ็นเซอร์แบบไม่ปรับเทียบ & สาธารณะ \\
\textcite{basavaraju2020} & สหรัฐอเมริกา & 3 ประเภท ได้แก่ ถนนเรียบ หลุม และรอยแตกตามขวางลึก & มาตรวัดความเร่ง มาตรวัดการหมุน ระบบระบุตำแหน่งบนพื้นโลก และวิดีโอสำหรับช่วยกำหนดป้ายกำกับ & ไม่สาธารณะ \\
\end{longtable}
\endgroup

\textcite{gonzalez2017} นำเสนอชุดข้อมูลสำคัญสำหรับงานจำแนกความผิดปกติของผิวถนน โดยเก็บข้อมูลจากมาตรวัดความเร่งของสมาร์ตโฟนในยานพาหนะ 12 คัน และครอบคลุมตำแหน่งการวางสมาร์ตโฟนหลายรูปแบบในเมืองชิวาวา ประเทศเม็กซิโก ชุดข้อมูลดังกล่าวประกอบด้วยเหตุการณ์ถนน 5 ประเภท ได้แก่ หลุม ลูกระนาดโลหะ ลูกระนาดยางมะตอย ถนนสึกหรอ และถนนปกติ งานนี้เสนอการแทนข้อมูลแบบ Bag-of-Words เพื่อแปลงสัญญาณมาตรวัดความเร่งแบบอนุกรมเวลาให้อยู่ในรูปเวกเตอร์คุณลักษณะขนาดคงที่สำหรับป้อนเข้าสู่ตัวจำแนก ผู้เขียนระบุว่างานวิจัยก่อนหน้าขาดชุดข้อมูลกลางสำหรับการเปรียบเทียบวิธีการ ซึ่งสะท้อนให้เห็นว่าการสร้างชุดข้อมูลภาคสนามมีคุณค่าต่อชุมชนวิจัย อย่างไรก็ตาม แม้บทความต้นฉบับจะระบุว่าชุดข้อมูลถูกเผยแพร่สาธารณะ แต่รายงานทบทวนล่าสุดระบุว่าชุดข้อมูลดังกล่าวไม่สามารถเข้าถึงได้แล้วในปัจจุบัน

เมื่อพิจารณาความเกี่ยวข้องกับงานวิจัยนี้ \textcite{gonzalez2017} มีจุดเด่นด้านการสร้างชุดข้อมูลอ้างอิงและการทดสอบวิธีจำแนกบนข้อมูลที่เก็บจากหลายยานพาหนะและหลายตำแหน่งการวางสมาร์ตโฟน อย่างไรก็ตาม งานดังกล่าวมุ่งเน้นสัญญาณมาตรวัดความเร่งเป็นหลัก โดยเฉพาะการใช้ข้อมูลแกนตั้งและการแทนข้อมูลแบบ Bag-of-Words ขณะที่งานวิจัยนี้มุ่งสร้างชุดข้อมูลภาคสนามหลายแหล่งในพื้นที่เฉพาะ โดยประกอบด้วยข้อมูลจากหน่วยวัดความเฉื่อย ระบบระบุตำแหน่งบนพื้นโลกและความเร็ว วิดีโอสำหรับตรวจสอบป้ายกำกับ คุณลักษณะทางเสียงที่ผ่านการประมวลผล และกระบวนการกำหนดป้ายกำกับภายใต้บริบทถนนเมืองนครราชสีมา ความแตกต่างสำคัญจึงไม่ได้อยู่ที่การเสนอชุดข้อมูลเพียงอย่างเดียว แต่อยู่ที่การออกแบบกระบวนการสร้างข้อมูลภาคสนามที่สามารถตรวจสอบและทำซ้ำได้ ตั้งแต่โปรแกรมประยุกต์ โครงสร้างฐานข้อมูล กระบวนการกำหนดป้ายกำกับ ไปจนถึงองค์ประกอบสำหรับการประเมินผล

\textcite{khandakar2025} นำเสนอชุดข้อมูลสมาร์ตโฟนแบบหลายเซ็นเซอร์สำหรับการวิเคราะห์ความปลอดภัยทางถนนและสภาพถนนในเมืองราชชาฮี ประเทศบังกลาเทศ ชุดข้อมูลนี้ครอบคลุมข้อมูลจากมาตรวัดความเร่ง มาตรวัดการหมุน สนามแม่เหล็ก ระบบระบุตำแหน่งบนพื้นโลก แรงโน้มถ่วง ทิศทาง ความเร่งรวม และข้อมูลจากเซ็นเซอร์แบบไม่ปรับเทียบ งานดังกล่าวมีความสำคัญในฐานะชุดข้อมูลสาธารณะสมัยใหม่ที่แสดงให้เห็นแนวโน้มของการใช้ข้อมูลหลายแหล่งจากสมาร์ตโฟนเพื่อสนับสนุนการวิเคราะห์สภาพถนนและพฤติกรรมการขับขี่ ผู้เขียนยังระบุว่าการขาดชุดข้อมูลมาตรฐานที่เชื่อถือได้เป็นข้อจำกัดสำคัญของงานวิจัยด้านนี้

อย่างไรก็ตาม แม้ \textcite{khandakar2025} จะเผยแพร่ชุดข้อมูลสาธารณะและครอบคลุมเซ็นเซอร์หลายชนิด แต่งานดังกล่าวยังอยู่ในบริบทพื้นที่ของประเทศบังกลาเทศ และมุ่งเน้นทั้งเหตุการณ์ถนนผิดปกติและพฤติกรรมการขับขี่เป็นหลัก ไม่ได้มุ่งสร้างชุดข้อมูลอ้างอิงเฉพาะบริบทถนนเมืองไทยหรือเมืองนครราชสีมา อีกทั้งยังไม่ได้เน้นกระบวนการกำหนดป้ายกำกับย้อนหลังระหว่างการเก็บข้อมูลภาคสนามร่วมกับวิดีโอและคุณลักษณะทางเสียงในลักษณะเดียวกับงานวิจัยนี้ ดังนั้น งานของ Khandakar et al. จึงช่วยยืนยันความสำคัญของการสร้างชุดข้อมูลสมาร์ตโฟนแบบหลายเซ็นเซอร์ แต่ยังคงเหลือช่องว่างด้านบริบทพื้นที่ กระบวนการกำหนดป้ายกำกับ และกรอบการประเมินที่สามารถทำซ้ำได้สำหรับงานวิจัยนี้

\textcite{basavaraju2020} ศึกษาการจำแนกสภาพถนนจากข้อมูลสมาร์ตโฟน โดยใช้ข้อมูลจากมาตรวัดความเร่ง มาตรวัดการหมุน และระบบระบุตำแหน่งบนพื้นโลก เพื่อจำแนกสภาพถนนออกเป็น 3 ประเภท ได้แก่ ถนนเรียบ หลุม และรอยแตกตามขวางลึก งานนี้มีจุดเด่นที่ใช้ข้อมูลจากหลายแกนของเซ็นเซอร์ และใช้วิดีโอช่วยสนับสนุนการกำหนดป้ายกำกับของสภาพถนน ผู้วิจัยพัฒนาโปรแกรมช่วยกำหนดป้ายกำกับจากวิดีโอ โดยสามารถตรวจสอบตำแหน่งที่ล้อรถผ่านเหตุการณ์ผิดปกติของถนนก่อนกำหนดป้ายกำกับให้กับข้อมูลสัญญาณ

อย่างไรก็ตาม \textcite{basavaraju2020} เป็นงานที่เก็บข้อมูลในบริบทของสหรัฐอเมริกา และไม่ได้เผยแพร่ชุดข้อมูลเป็นสาธารณะสำหรับใช้เป็นข้อมูลอ้างอิงกลาง นอกจากนี้ แม้งานดังกล่าวจะใช้วิดีโอช่วยกำหนดป้ายกำกับ แต่ยังไม่ได้มุ่งออกแบบกระบวนการกำหนดป้ายกำกับย้อนหลังสำหรับการใช้งานภาคสนามจริงในลักษณะที่รองรับการกดป้ายกำกับระหว่างการขับขี่ การตรวจสอบภายหลัง และการคัดแยกเหตุการณ์ที่ไม่สามารถตัดสินได้อย่างชัดเจนออกจากชุดข้อมูลฝึกสอน งานนี้จึงมีความเกี่ยวข้องในฐานะตัวอย่างของการใช้สมาร์ตโฟนและวิดีโอเพื่อช่วยสร้างป้ายกำกับ แต่ยังไม่ครอบคลุมกรอบการสร้างชุดข้อมูลอ้างอิงที่เปิดเผยกระบวนการและสามารถทำซ้ำได้ครบถ้วน

\textcite{silva2018} แสดงให้เห็นว่าระบบตรวจจับความผิดปกติของถนนที่ประเมินภายใต้สภาวะควบคุมอาจมีประสิทธิภาพลดลงเมื่อนำไปใช้กับสภาพแวดล้อมถนนจริง ผลการศึกษาดังกล่าวสนับสนุนว่า ข้อมูลภาคสนามจริงไม่ได้เป็นเพียงข้อมูลประกอบ แต่เป็นเงื่อนไขสำคัญต่อความสามารถในการใช้งานทั่วไปและความถูกต้องของแบบจำลองเมื่อนำไปประยุกต์ใช้จริง อย่างไรก็ตาม วัตถุประสงค์ของ \textcite{silva2018} มุ่งเน้นการประเมินระบบที่พัฒนาแล้วในสภาพการใช้งานจริง ขณะที่งานวิจัยนี้มุ่งสร้างชุดข้อมูลอ้างอิงและกรอบการประเมินที่สามารถทำซ้ำได้ ซึ่งประกอบด้วยชุดข้อมูลอ้างอิงที่ผ่านการตรวจสอบภาคสนาม ระเบียบวิธีการกำหนดป้ายกำกับที่ชัดเจน รหัสโปรแกรม การตั้งค่า และการแบ่งชุดข้อมูลสำหรับตรวจสอบและทำซ้ำผลการทดลอง ความแตกต่างด้านระดับความละเอียดของความสามารถในการทำซ้ำนี้ทำให้งานของ Silva et al. ไม่สามารถใช้เป็นฐานสำหรับการเปรียบเทียบแบบจำลองหลายรูปแบบภายใต้กระบวนการประเมินเดียวกันได้โดยตรง

\textcite{chibani2022} นำเสนองานประเมินวิธีการแบ่งช่วงสัญญาณ โดยเปรียบเทียบการแบ่งช่วงสัญญาณแบบเลื่อนหน้าต่างพลวัตกับการแบ่งช่วงสัญญาณแบบเลื่อนหน้าต่างคงที่ สำหรับงานตรวจจับความผิดปกติของถนนจากสัญญาณมาตรวัดความเร่ง งานดังกล่าวทดสอบบนชุดข้อมูล Pothole Lab และชุดข้อมูล CarSim ผลการทดลองแสดงให้เห็นว่าการแบ่งช่วงสัญญาณแบบเลื่อนหน้าต่างพลวัตให้ผลดีกว่าการแบ่งช่วงสัญญาณแบบเลื่อนหน้าต่างคงที่อย่างสม่ำเสมอ ทั้งกับตัวตรวจจับเชิงกฎและตัวจำแนกแบบการเรียนรู้ของเครื่องหลายรูปแบบ เช่น KNN, RF และ SVM ผลลัพธ์ดังกล่าวสนับสนุนแนวคิดว่า การแบ่งช่วงสัญญาณตามความเร็วมีเหตุผลในเชิงกายภาพ เนื่องจากความยาวของสัญญาณที่เกิดจากเหตุการณ์ผิดปกติบนถนนเปลี่ยนแปลงตามความเร็วของยานพาหนะ ไม่ได้มีขนาดคงที่เสมอ

งานวิจัยนี้จึงนำการแบ่งช่วงสัญญาณแบบขนาดคงที่และการแบ่งช่วงสัญญาณแบบปรับตามความเร็วมาเปรียบเทียบกันเป็นส่วนหนึ่งของการประเมินบนข้อมูลภาคสนามจริง จุดเน้นของงานวิจัยนี้ไม่ใช่การเสนอการแบ่งช่วงสัญญาณแบบปรับตามความเร็วเป็นวิธีใหม่ แต่เป็นการประเมินแนวคิดดังกล่าวภายใต้ข้อมูลถนนเมืองจริงที่มีความแปรปรวนของความเร็ว พฤติกรรมการจราจร ความแตกต่างของตำแหน่งการติดตั้งสมาร์ตโฟน และกระบวนการกำหนดป้ายกำกับที่ตรวจสอบได้ การประเมินในลักษณะนี้ช่วยให้สามารถพิจารณาได้ว่าแนวคิดที่ให้ผลดีในชุดข้อมูลจำลองหรือชุดข้อมูลอ้างอิงเดิม สามารถถ่ายโอนมาใช้กับบริบทภาคสนามจริงของเมืองนครราชสีมาได้มากน้อยเพียงใด

จากการทบทวนงานวิจัยที่เกี่ยวข้อง พบว่างานวิจัยด้านการตรวจสอบสภาพถนนด้วยสมาร์ตโฟนมีความก้าวหน้าในหลายด้าน ทั้งการสร้างชุดข้อมูล การใช้ข้อมูลหลายเซ็นเซอร์ การใช้วิดีโอช่วยกำหนดป้ายกำกับ และการศึกษาวิธีแบ่งช่วงสัญญาณตามความเร็ว อย่างไรก็ตาม ยังมีช่องว่างสำคัญหลายประการ ได้แก่ การขาดชุดข้อมูลอ้างอิงที่ผ่านการตรวจสอบภาคสนามสำหรับถนนในบริบทไทย การขาดระเบียบวิธีการกำหนดป้ายกำกับที่เหมาะสมกับการเก็บข้อมูลระหว่างการขับขี่จริง การขาดชุดข้อมูลที่รวมข้อมูลจากหน่วยวัดความเฉื่อย ระบบระบุตำแหน่ง วิดีโอ และคุณลักษณะทางเสียงที่ผ่านการประมวลผลในกระบวนการเดียวกัน และการขาดกรอบการประเมินที่สามารถทำซ้ำได้สำหรับการเปรียบเทียบวิธีจำแนกหลายรูปแบบบนชุดข้อมูลเดียวกัน

ดังนั้น งานวิจัยนี้จึงมุ่งเติมช่องว่างดังกล่าวโดยการพัฒนาแพลตฟอร์มและกระบวนการสร้างชุดข้อมูลภาคสนามสำหรับการจำแนกสภาพถนนในเขตเมืองนครราชสีมา ชุดข้อมูลที่สร้างขึ้นจะประกอบด้วยข้อมูลจากหน่วยวัดความเฉื่อย ระบบระบุตำแหน่งบนพื้นโลกและความเร็ว วิดีโอสำหรับตรวจสอบป้ายกำกับ คุณลักษณะทางเสียงที่ผ่านการประมวลผล และป้ายกำกับที่กำหนดผ่านกระบวนการย้อนหลังและตรวจสอบได้ นอกจากนี้ งานวิจัยยังจัดเตรียมองค์ประกอบด้านความสามารถในการทำซ้ำ ได้แก่ โปรแกรมประยุกต์ โครงสร้างฐานข้อมูล กระบวนการเตรียมข้อมูล การแบ่งชุดข้อมูล การตั้งค่าการทดลอง และกรอบการประเมิน เพื่อให้สามารถใช้เป็นฐานสำหรับการเปรียบเทียบวิธีจำแนกสภาพถนนภายใต้บริบทการใช้งานจริงได้อย่างเป็นระบบมากขึ้น.
